\subsection{Temporal smoothing}
\label{sec:method-smoothing}

Sleep stages are strongly autocorrelated: an epoch of N3 is far more likely
to be followed by N3 than by W, and a per-epoch softmax discards that
information. We fit a first-order hidden Markov model to the training labels,
with transition matrix $A$ (Laplace-smoothed) and initial distribution $\pi$,
and run the forward--backward recursions in log space \emph{within each
recording} to obtain
\begin{equation}
\begin{aligned}
\gamma_t(k)&=\frac{a_t(k)b_t(k)}{\sum_{k'\in\cY}a_t(k')b_t(k')},\\
a_1(k)&=\pi_k\,p_1(k),\qquad
a_t(k)=p_t(k)\sum_{j}a_{t-1}(j)A_{jk}\quad(t\ge2),
\end{aligned}
\label{eq:forward-backward}
\end{equation}
and $b_T(k)=1$,
$b_t(k)=\sum_{j}A_{kj}\,p_{t+1}(j)\,b_{t+1}(j)$. (The forward and backward messages are
written $a$ and $b$ so that $\alpha$ and $\beta$ keep their meanings elsewhere
in the paper: the nominal level and its running ACI target, and the block
transition weight.) The recursion is $O(LK^2)$ per recording.

The smoothing uses no test labels and no test-time adaptation: it is a
deterministic transformation of the model's own output. If calibration and
test epochs are exchangeable under the score $s(\cdot,\cdot)$, they remain
exchangeable under $s(\gamma(\cdot),\cdot)$, so split CP stays valid and the
change is one of \emph{efficiency}, not of validity. That argument is the
reason we prefer smoothing to a post-hoc repair of the sets, which would
break the exchangeability argument it relies on.

\subsection{Pooled calibration and its failure mode}
\label{sec:method-pooled}

Pooled calibration guarantees that $1-\alpha$ of \emph{all} epochs are
covered, not that $1-\alpha$ of a given patient's epochs are. With $C$ the
set of calibration subjects and $S_{j,t}$ the score of epoch $t$ of subject
$j$,
\begin{equation}
\wh q_g=\operatorname{Quantile}_{\lceil(n_C+1)(1-\alpha)\rceil/n_C}
\big(\{S_{j,t}\}_{j\in C}\big),
\label{eq:global-q}
\end{equation}
which is \Cref{eq:split-cp} applied to the pooled scores. \Cref{sec:results-patient}
quantifies how far per-patient coverage departs from nominal under this rule.

\subsection{Subject-adaptive calibration with a warm-up}
\label{sec:method-sacp}

If a new subject supplies $m$ labelled epochs before deployment, those epochs
define a local threshold $\wh q_i$ by the same rule, and we deploy
\begin{equation}
\wh q_i(\lambda)=(1-\lambda)\wh q_g+\lambda\wh q_i,
\qquad \lambda\in[0,1].
\label{eq:shrunk-q}
\end{equation}
The shrinkage matters because a warm-up is short: with $m\approx 1/\alpha$
the local threshold is an extreme order statistic and its sampling variance
is large, so $\lambda=1$ trades a bias reduction for a variance increase.

\begin{proposition}[Shrinkage trade-off]
\label{prop:shrinkage}
Let subject $i$ have score CDF $F_i$ that is $\ell$-Lipschitz in a neighbourhood
of its $(1-\alpha)$-quantile $q_i^\star$, let $\Delta=q_i^\star-q_g^\star$ be
the subject's offset from the cohort quantile, and suppose $\wh q_i$ and
$\wh q_g$ are approximately unbiased with variances $\sigma^2/m$ and
$\sigma^2/n_C$. Then for a threshold $\wh q_i(\lambda)$ as in
\Cref{eq:shrunk-q}, the squared coverage gap satisfies
\begin{equation}
\bE\Big[\big(F_i(\wh q_i(\lambda))-(1-\alpha)\big)^2\Big]
\ \le\ \ell^2\Big[(1-\lambda)^2\Delta^2+\lambda^2\frac{\sigma^2}{m}
 +(1-\lambda)^2\frac{\sigma^2}{n_C}\Big],
\label{eq:shrink-bound}
\end{equation}
whose minimiser over $\lambda$ is
\begin{equation}
\lambda^\star=\frac{\Delta^2+\sigma^2/n_C}{\Delta^2+\sigma^2/n_C+\sigma^2/m}.
\label{eq:lambda-star}
\end{equation}
\end{proposition}

\begin{proof}[Proof sketch]
Write $F_i(\wh q_i(\lambda))-(1-\alpha)=F_i(\wh q_i(\lambda))-F_i(q_i^\star)$,
which is bounded in absolute value by $L|\wh q_i(\lambda)-q_i^\star|$ by the
Lipschitz condition. Expanding $\wh q_i(\lambda)-q_i^\star
=(1-\lambda)(\wh q_g-q_g^\star-\Delta)+\lambda(\wh q_i-q_i^\star)$ and taking
expectations of squares under independence of the two estimates gives
\Cref{eq:shrink-bound}; differentiating with respect to $\lambda$ and setting
the derivative to zero gives \Cref{eq:lambda-star}.
\end{proof}

Proposition~\ref{prop:shrinkage} is not a plug-in estimator: $L$, $\sigma$ and
$\Delta$ are unknown at deployment, and the Gaussian approximation to quantile
estimates is asymptotic in $m$. We use it for what it does predict --- that
$\lambda^\star\to0$ as $m\to0$, that $\lambda^\star\to1$ as $m\to\infty$, and
that adaptation pays only when the subject's offset dominates the calibration
noise, $\Delta^2\gg\sigma^2/n_C$ --- and we select $\lambda$ by leave-one-subject-out
simulation on the calibration subjects, scoring each candidate by the
one-sided mean shortfall
\begin{equation}
\mathcal{L}(\lambda)=\bE\big[\max\{0,(1-\alpha)-\widehat{\mathrm{cov}}\}\big],
\label{eq:shortfall}
\end{equation}
so that over-coverage is not rewarded; it is paid for in set size and is not
banked as a gain. Test subjects never enter this selection, which differs
from the usual practice of reporting the best $\lambda$ on test data.

\paragraph{Label-free variant}
When no warm-up labels exist, each subject is described by unlabelled
statistics $z_i=(\bar p_i,\ \overline{\max_k p_{i,k}},\ \bar H(p_i))\in\bR^{7}$
--- the mean predicted class distribution, the mean confidence and the mean
predictive entropy, with the distribution contributing five coordinates ---
and the $k$ nearest calibration subjects in this space supply a local
threshold, shrunk toward the global one with the same $\lambda$. This is an
approximation and we report it as such: it cannot be valid in general, only
when the descriptor correlates with the subject's score distribution.

\subsection{Per-patient adaptive conformal inference}
\label{sec:method-aci}

The warm-up regime presumes labelled epochs from the new patient before
deployment. In a review workflow the opposite is true: labels arrive
\emph{after} each prediction, because a technologist reviews the automatic
hypnogram. This is the feedback pattern ACI was designed for.

Running a single ACI stream over all epochs of all patients gives a pooled
guarantee once more: a hard patient can be under-covered while an easy one is
over-covered and the stream still converges. We therefore run one
\emph{independent} stream per patient $p$,
\begin{equation}
\alpha_{p,t+1}=\alpha_{p,t}+\gamma_{p,t}\big(\alpha-\ones\{Y_{p,t}\notin\cC_{p,t}\}\big),
\qquad
\cC_{p,t}=\{y:\ s_{p,t}(y)\le\wh q(\alpha_{p,t})\},
\label{eq:aci-patient}
\end{equation}
each starting at $\alpha_{p,1}=\alpha$ and each reading its threshold from the
same pooled calibration scores, so that no patient-specific calibration data
is required. The step schedule is
\begin{equation}
\gamma_{p,t}=\gamma_0/\sqrt{t+2},
\label{eq:step}
\end{equation}
which decays fast enough to stop a single misclassified stretch from
permanently loosening the threshold, and slowly enough to keep adapting over a
night. The schedule has $\sum_t\gamma_t=\infty$ but also
$\sum_t\gamma_t^2=\infty$, so it violates the square-summability condition
under which \citet{angelopoulos2024} establish their guarantee; we make
\emph{no} long-run guarantee for it, and report
the realised per-patient coverage and state in \Cref{sec:results-patient}
whether it is above or below nominal. What the procedure buys, and what we
measure, is the reallocation of coverage from easy to hard patients.

The point of running one stream per patient is that it turns ACI's pooled
long-run statement into a per-patient one. We state that implication
explicitly, because it is the only part of the construction that carries a
guarantee.

\begin{corollary}[Per-patient long-run coverage]
\label{cor:per-patient}
Suppose the threshold map $\wh q(\alpha)$ is the fixed quantile of the pooled
calibration scores, so that \Cref{eq:aci-patient} defines every set, and let
the step schedule satisfy $\gamma_t\to0$, $\gamma_t\le1$,
$\sum_t\gamma_t=\infty$ and $\sum_t\gamma_t^2<\infty$. Then for every patient
$p$ the realised miscoverage rate on that patient's stream converges to
$\alpha$,
\begin{equation}
\frac{1}{T}\sum_{t=1}^{T}\ones\{Y_{p,t}\notin\cC_{p,t}\}
\;\longrightarrow\;\alpha
\qquad (T\to\infty),
\label{eq:aci-per-patient}
\end{equation}
for an arbitrary sequence of inputs and labels, with no exchangeability
assumption within the patient.
\end{corollary}

\begin{proof}
The update \Cref{eq:aci-patient} of stream $p$ depends only on that stream's
own outcomes and on a threshold map fixed by the calibration set, so the
stream is an instance of the single-stream recursion of
\citet{angelopoulos2024}. Their long-run bound, applied stream by stream,
gives \Cref{eq:aci-per-patient}. The step \Cref{eq:step} used in the
experiments does not satisfy the summability hypothesis and therefore carries
no claim; \Cref{sec:results-aci} repeats the experiment with a schedule that
does.
\end{proof}

The scale $\gamma_0$ is a hyperparameter. We choose it by leave-one-subject-out
simulation within the calibration cohort: each calibration subject in turn
forms a stream, the remaining calibration subjects supply the score pool, and
the $\gamma_0$ that minimises the one-sided mean shortfall
\Cref{eq:shortfall} across those held-out streams --- ties broken by smaller
mean set size --- is applied unchanged to every test patient. The test cohort
never influences the choice.

\subsection{Sequence-level constructors}
\label{sec:method-seq}

A per-epoch set is free to be inconsistent with its neighbours: epoch $t$ may
admit $\{W,N1\}$ while $t+1$ admits $\{N3\}$, a transition the training labels
essentially forbid. We derived two constructors that repair this by scoring
labels in the context of the best \emph{sequence} through them, and we report
both with their properties because one of them is a useful ablation and the
other is a negative result.

\paragraph{SeqCP-Margin}
For a block of $L$ consecutive epochs inside one recording, transition cost
$\tau(a,b)=1-A_{ab}$, and weight $\beta\ge0$, define
\begin{equation}
\begin{aligned}
M_t(y)&=\min_{y_{1:L}:\,y_t=y}
\Big[\frac1L\sum_{u=1}^{L}\big(1-\gamma_u(y_u)\big)
+\frac{\beta}{L-1}\sum_{u=2}^{L}\tau(y_{u-1},y_u)\Big],\\
M^\star&=\min_{t,y}M_t(y).
\end{aligned}
\label{eq:margin}
\end{equation}
Forward and backward recursions compute every $M_t(y)$ in $O(LK^2)$; a
brute-force check over all $K^L$ paths for small blocks is part of the test
suite. We take $m_t=M_t(Y_t)-M^\star$ as the nonconformity of the true label,
calibrate $\wh\delta$ as its split-conformal quantile on calibration blocks,
and predict $\cC_t=\{y:M_t(y)-M^\star\le\wh\delta\}$. Since
$Y_t\in\cC_t\iff m_t\le\wh\delta$, this is split CP on the scalar $m_t$ and
inherits its marginal guarantee.

\begin{proposition}[Reduction at $\beta=0$]
\label{prop:reduction}
With $\beta=0$ the dynamic program factorises: for every $t$ and $y$,
\begin{equation}
M_t(y)-M^\star=\frac{1}{L}\Big(s_t(y)-\min_{k\in\cY}s_t(k)\Big),
\qquad s_t(k)=1-\gamma_t(k),
\label{eq:reduction}
\end{equation}
independently of the block. Hence SeqCP-Margin at $\beta=0$ is exactly
per-epoch split CP on the deficit score \Cref{eq:deficit}, and any difference
between the two at $\beta>0$ is attributable to the transition term alone.
\end{proposition}

\begin{proof}
At $\beta=0$ the objective separates over positions. The minimum total cost of
a path constrained to $y_t=y$ is $s_t(y)+\sum_{u\ne t}\min_k s_u(k)$, so
$M_t(y)-M^\star=\frac1L(s_t(y)-\min_k s_t(k))$ after subtracting the
unconstrained minimum $\frac1L\sum_u\min_k s_u(k)$.
\end{proof}

\begin{proposition}[Monotonicity of the centring]
\label{prop:monotone}
Let $M_t^{\mathrm{pos}}(y)=M_t(y)-\min_k M_t(k)$. Then
$M_t^{\mathrm{pos}}(y)\le M_t(y)-M^\star$ for every $t,y$, so the
position-centred sets are pointwise no larger than the globally centred ones.
\end{proposition}

\begin{proof}
$\min_k M_t(k)\ge\min_{u,k}M_u(k)=M^\star$.
\end{proof}

In practice the two coincide, because the block-optimal path on a sleep block
is usually constant and therefore visits every epoch, in which case
$\min_k M_t(k)=M^\star$ for all $t$. We keep both and report that they agree
rather than asserting a difference.

\paragraph{SeqCP-Block}
Projecting a \emph{sequence-level} conformal set onto each position gives a
simultaneous guarantee with no union bound. Calibrate the block score
\begin{equation}
S(y_{1:L})=\frac1L\sum_{t=1}^{L}\aps_t(y_t)
+\frac{\beta}{L-1}\sum_{t=2}^{L}\tau(y_{t-1},y_t)
\label{eq:block-score}
\end{equation}
on complete calibration blocks, obtaining $\wh q$, and predict
$\cC_t=\{y:M_t(y)\le\wh q\}$ with $M_t$ computed from the APS score matrix.

\begin{theorem}[Simultaneous coverage of the projection]
\label{thm:projection}
Let $\wh q$ be the split-conformal quantile of \Cref{eq:block-score} over
$n_B$ exchangeable calibration blocks. Then with probability at least
$1-\alpha$, the true block is in the block-level set, and in that
event $Y_t\in\cC_t$ for every $t=1,\dots,L$.
\end{theorem}

\begin{proof}
By \Cref{eq:split-cp}, $\bP(S(Y_{1:L})\le\wh q)\ge 1-\alpha$. If the true
block is in the block-level set then the path $y_{1:L}=Y_{1:L}$ is feasible,
so $M_t(Y_t)\le S(Y_{1:L})\le\wh q$ by the definition of $M_t$ as a minimum
over paths. Hence $Y_t\in\cC_t$ for all $t$ simultaneously.
\end{proof}

\Cref{thm:projection} is the guarantee we can prove, and
\Cref{sec:results-temporal} reports what it costs. The projection of a union
of sequences contains every label that \emph{any} admissible sequence assigns
to $t$, so it is inherently loose, and because the transition terms are
non-negative penalties the block set prefers low-transition sequences. At
$\beta\ge0.5$ the penalised block set admits a constant sequence for every
label, the projection becomes the whole label set, and the guarantee is
satisfied trivially. Even at $\beta=0$ the block mean lets one well-scored
neighbour absorb a badly scored epoch. We report this as a property of the
construction rather than of our implementation.

\subsection{What we do not claim}
\label{sec:method-claims}

The nonconformity scores are calibrated, not the model: a set can cover the
truth and still be clinically useless if it is always $\{N1,N2\}$, which is
why set size, singleton rate and review workload are reported alongside
coverage. We make no subject-conditional finite-sample guarantee under
arbitrary shift, no claim that SeqCP-Block is efficient, and no clinical
safety claim. The assumptions each result needs are stated where the result
is used.
